\chapter{OUTPUT STRUCTURE AND RELEASE WORKFLOW}

\noindent\textbf{Mục tiêu chương}

\begin{enumerate}[label=-]
\item Mô tả cấu trúc output hiện hành của hệ thống báo cáo năng lượng.
\item Ghi nhận release workflow, naming conventions, và các manual smoke checks đang được chấp nhận.
\item Tóm tắt PDF staging, print-path baseline, và release-readiness rules đang được dùng thật.
\item Giữ chapter này tập trung vào output artifacts và release workflow, còn deployment và troubleshooting được tách sang chapter riêng.
\end{enumerate}

\section{Vai trò của chapter output/release}

Sau khi đã mô tả kiến trúc, data flow, và business rules, bộ tài liệu kỹ thuật cần có một chapter output/release để trả lời các câu hỏi thực tế sau:

\begin{enumerate}[label=-]
\item hệ thống đang ghi output ở đâu;
\item cần kiểm tra gì trước khi gọi một bản build là ổn để phát hành;
\item manual smoke anchors nào đang được chấp nhận cho runtime hiện tại;
\item khi nào documentation PDF cần rebuild sau một thay đổi ở source docs.
\end{enumerate}

Chapter này được tổng hợp từ workflow thực tế của dự án và nên được giữ đồng bộ với các tài liệu nguồn như \texttt{docs/workflows/release\_runbook.md}, \texttt{docs/workflows/release\_checklist.md}, và \texttt{docs/workflows/pdf\_export\_flow.md}. Phần bootstrap host mới và schedule enable flow được tách sang chapter \textit{Deployment and Operations} để giảm việc trộn lẫn giữa release rules và operator setup.

\section{Canonical output structure}

Hệ thống hiện ghi artifacts theo tháng, thay vì dùng dated run folders như trước. Root chuẩn là \texttt{output/reports/YYYY\_MM/}. Mỗi tháng được tách thành các nhóm output sau:

\begin{enumerate}[label=-]
\item \texttt{view\_html/}: HTML view phục vụ đọc và kiểm tra trên browser;
\item \texttt{pdf\_source\_html/}: HTML nguồn dùng cho PDF export và kiểm tra print layout;
\item \texttt{pdf/}: PDF cuối cùng sau khi Chromium hoặc Chrome DevTools Protocol hoàn tất print;
\item \texttt{excel/}: workbook Excel daily-only của version hiện tại.
\end{enumerate}

Bảng \ref{tab:artifact_structure} tóm tắt output conventions đang được chấp nhận.

\begin{table}[h]
	\centering
	\renewcommand{\arraystretch}{1.3}
	\caption{Canonical output artifact structure}
	\label{tab:artifact_structure}
	\begin{tabularx}{\linewidth}{|>{\RaggedRight\arraybackslash}p{2.6cm}|>{\RaggedRight\arraybackslash}p{3.2cm}|>{\RaggedRight\arraybackslash}X|}
		\hline
		\textbf{Artifact Family} & \textbf{Location} & \textbf{Purpose} \\ \hline
		View HTML & \TablePath{output/reports/YYYY_MM/view_html/} & Browser-readable version dùng để đọc report, kiểm tra cấu trúc, và debug render output. \\ \hline
		PDF source HTML & \TablePath{output/reports/YYYY_MM/pdf_source_html/} & HTML nguồn cho PDF export, phản ánh print-safe shell và compact print styling trước khi in. \\ \hline
		Final PDF & \TablePath{output/reports/YYYY_MM/pdf/} & PDF A4 cuối cùng, dùng cho phát hành và lưu trữ. \\ \hline
		Daily Excel workbook & \TablePath{output/reports/YYYY_MM/excel/} & Workbook daily-only của v1 Excel export, phục vụ tabular extraction. \\ \hline
	\end{tabularx}
\end{table}

\subsection{Filename ordering rules}

Filename ordering được kiểm soát bằng prefix số để việc sort nhất quán trên nhiều môi trường file system:

\begin{enumerate}[label=-]
\item \texttt{01\_monthly}
\item \texttt{02\_weekly}
\item \texttt{03\_daily}
\end{enumerate}

Base name hiện tại của report artifacts là \texttt{energy\_automatic\_report}. Như vậy các ví dụ canonical gồm:

\begin{enumerate}[label=-]
\item \texttt{01\_monthly\_energy\_automatic\_report\_20250531.pdf}
\item \texttt{02\_weekly\_energy\_automatic\_report\_20250629.pdf}
\item \texttt{03\_daily\_energy\_automatic\_report\_20250629.xlsx}
\end{enumerate}

Naming discipline này quan trọng vì nó ảnh hưởng trực tiếp đến khả năng rà soát output, diff artifact giữa các lần chạy, và logic release smoke checks.

\section{Runtime scheduling và artifact expectations}

Runtime hiện tại tuân theo các scheduling rules sau:

\begin{enumerate}[label=-]
\item luôn export \textit{daily};
\item export \textit{weekly} khi anchor day là Chủ nhật;
\item export \textit{monthly} khi anchor day là ngày cuối tháng;
\item nếu cả hai điều kiện cùng đúng, runtime export đầy đủ các report tương ứng trong một lần chạy.
\end{enumerate}

Bảng \ref{tab:scheduled_artifacts} tóm tắt kỳ vọng artifacts theo từng tình huống smoke test tiêu biểu.

\begin{table}[h]
	\centering
	\renewcommand{\arraystretch}{1.3}
	\caption{Expected artifacts for common smoke scenarios}
	\label{tab:scheduled_artifacts}
	\begin{tabularx}{\linewidth}{|>{\RaggedRight\arraybackslash}p{2.7cm}|>{\RaggedRight\arraybackslash}p{2.6cm}|>{\RaggedRight\arraybackslash}X|}
		\hline
		\textbf{Scenario} & \textbf{Expected Periods} & \textbf{Expected Canonical Outputs} \\ \hline
		Daily-only anchor & Daily & view HTML, PDF source HTML, final PDF, và daily Excel workbook dưới prefix \texttt{03\_daily}. \\ \hline
		Sunday anchor & Daily + Weekly & toàn bộ artifact daily cộng thêm bộ weekly dưới prefix \texttt{02\_weekly}; không có monthly. \\ \hline
		Month-end anchor & Daily + Monthly & toàn bộ artifact daily cộng thêm bộ monthly dưới prefix \texttt{01\_monthly}; không có weekly. \\ \hline
		Sunday + month-end & Daily + Weekly + Monthly & runtime xuất đủ cả ba period families trong cùng tháng output tương ứng. \\ \hline
	\end{tabularx}
\end{table}

\section{PDF staging và print path}

PDF export của dự án tách rõ giữa canonical artifact path và staging path dùng cho Chromium print. Điều này giúp tránh các vấn đề với hidden directories và giữ cho output cuối cùng ổn định hơn.

\subsection{Canonical path và staging path}

\begin{enumerate}[label=-]
\item canonical project output root là \texttt{output/reports/};
\item mỗi report được ghi vào month-grouped directories bên dưới canonical root;
\item PDF print flow đồng thời tạo staging HTML và staging PDF ở một \textit{Chromium-safe non-hidden directory};
\item sau khi print thành công, staged PDF được copy ngược lại vào canonical PDF directory của tháng tương ứng.
\end{enumerate}

\subsection{Staging resolution order}

Thứ tự resolve staging directory hiện hành là:

\begin{enumerate}[label=-]
\item dùng \texttt{PRINT\_STAGING\_DIR} nếu được cấu hình rõ ràng;
\item nếu không có, dùng \texttt{OUTPUT\_DIR} khi path này không hidden;
\item nếu vẫn chưa có, dùng canonical project output dir nếu bản thân path đó không hidden;
\item fallback cuối cùng là \texttt{\textasciitilde/Reports}.
\end{enumerate}

Cách phân tầng này giúp tách bài toán artifact ownership với bài toán trình duyệt headless có thể truy cập được thư mục staging hay không.

\subsection{Stable PDF export baseline}

Print baseline hiện tại nên được xem là rule vận hành chính thức:

\begin{enumerate}[label=-]
\item ưu tiên Chrome DevTools Protocol print path;
\item dùng \texttt{scale=1.0};
\item dùng \texttt{preferCSSPageSize=true};
\item dùng \texttt{printBackground=true};
\item chỉ fallback sang Chromium CLI path khi CDP print thực sự thất bại.
\end{enumerate}

Đây là điểm mấu chốt để giữ sự ổn định giữa các report family và tránh hiện tượng document shrink hoặc auto-fit không mong muốn.

\section{Release workflow và smoke checks}

\subsection{Release workflow tối thiểu}

Trước khi một bản chạy được xem là đủ tốt cho release review, nên đi theo chu trình sau:

\begin{enumerate}[label=-]
\item xác nhận working copy và runtime environment đúng project path;
\item xác nhận Chrome hoặc Chromium hiện diện trên host;
\item kiểm tra các cấu hình như \texttt{REPORT\_ANCHOR\_DATE}, \texttt{OUTPUT\_DIR}, và \texttt{PRINT\_STAGING\_DIR};
\item chạy automated checks hoặc pytest phù hợp nếu chúng nằm trong scope hiện tại;
\item chạy manual smoke cho ít nhất một daily anchor;
\item chạy thêm Sunday anchor và month-end anchor để xác thực weekly/monthly scheduling khi cần;
\item mở PDF cuối để kiểm tra layout width, chart sizing, table overflow, và artifact copy-back;
\item ghi lại commit SHA và kết quả kiểm tra trước khi xin sign-off.
\end{enumerate}

\subsection{Manual smoke anchors}

Các anchor phổ biến đang được dùng trong runbook gồm:

\begin{enumerate}[label=-]
\item daily-only smoke: \texttt{2025-06-25};
\item Sunday smoke: \texttt{2025-06-29};
\item month-end smoke: \texttt{2025-05-31}.
\end{enumerate}

Các anchor này đủ để xác minh cả period scheduling, canonical naming, month grouping, và daily-only Excel behavior trong runtime hiện tại.

\section{Documentation PDF rebuild rule}

Subsystem tài liệu kỹ thuật có workflow riêng. PDF của bộ documentation không cần rebuild cho mọi thay đổi code. Chỉ nên rebuild khi một trong các nhóm sau thay đổi:

\begin{enumerate}[label=-]
\item source LaTeX hoặc Markdown của project documentation;
\item figures hoặc diagrams được chèn vào technical PDF;
\item bibliography hoặc references;
\item workflow rules hoặc release rules được mô tả trong technical manual.
\end{enumerate}

Quy tắc này giúp tài liệu luôn sát implementation nhưng không làm build pipeline bị ồn vì những thay đổi code không ảnh hưởng đến nội dung docs.

\section{Release-readiness checklist}

Trước khi xem một batch output hoặc một technical PDF là sẵn sàng cho bước review tiếp theo, nên xác nhận tối thiểu các điểm sau:

\begin{enumerate}[label=-]
\item naming và month grouping của artifacts đúng theo conventions hiện tại;
\item daily, weekly, monthly scheduling hoạt động đúng với anchor day tương ứng;
\item PDF cuối cùng không có lỗi scale, missing charts, hoặc overflow nghiêm trọng;
\item daily Excel workbook chỉ xuất hiện ở đúng daily scope;
\item docs, runbook, checklist, và technical PDF vẫn phản ánh runtime thật;
\item commit SHA và kết quả smoke test đã được ghi nhận trước khi xin sign-off.
\end{enumerate}

\section{Tổng kết chương}

Chapter này giữ riêng phần output artifacts, release workflow, PDF rebuild rules, và release-readiness checklist. Deployment setup và troubleshooting playbook đã được tách sang các chapter kế tiếp để người đọc tra cứu nhanh hơn trong tình huống handover hoặc vận hành thật.
